% Part: first-order-logic
% Chapter: model-theory
% Section: overspill

\documentclass[../../../include/open-logic-section]{subfiles}

\begin{document}

\olfileid{mod}{bas}{ove}
\olsection{上溢}

\begin{thm}
\ollabel{overspill} 若语句集合 $\Gamma$ 具有任意大的有限模型，则它也具有一个无穷模型。
\end{thm}

\begin{proof}
将 $\Gamma$ 的语言扩张，添加可数无穷多个新常项 $c_0$, $c_1$, \dots，并考虑集合 $\Gamma \cup \{c_i \neq c_j :
i \neq j\}$。$\Gamma$ 具有任意大的有限模型，是指对任意 $m >0$，存在 $n\ge m$ 使得 $\Gamma$ 有一个势为~$n$ 的模型。这意味着 $\Gamma \cup \{c_i \neq
c_j : i \neq j\}$ 是有限可满足的。根据紧致性，$\Gamma
\cup \{c_i \neq c_j : i \neq j\}$ 有一个模型 $\Struct M$，其论域必为无穷，因为它满足所有不等式 $c_i \neq c_j$。
\end{proof}

\begin{prop}
\ollabel{inf-not-fo}
不存在一阶语言的语句 $!A$，使得它在结构~$\Struct M$ 中为真当且仅当该结构的论域 $\Domain{M}$ 是无穷的。
\end{prop}

\begin{proof}
如果存在这样的 $!A$，其否定 $\lnot !A$ 将在且仅在所有有限结构中为真，因此它将具有任意大的有限模型但不具有无穷模型，这与 \olref{overspill} 矛盾。
\end{proof}

\end{document}